\subsection*{\texttt{print!} and \texttt{println!}: compile-time inspection}

The bang variants (\texttt{print!}, \texttt{println!}) run at \textbf{compile time}, not at runtime. They print to the compiler's output and are useful for inspecting shapes, extents, and type information without launching a kernel:

\begin{lstlisting}[language=croqtile,caption={A compile-time shape print}]
__co__ void check_shapes(f32 [3, 2] b) {
  print!("shape of b: ");
  println!("b.span = ", b.span);
}
\end{lstlisting}

When you compile this file, the compiler emits:

\begin{lstlisting}[language=croqtile,caption={The printed shape}]
shape of b: b.span = [3, 2]
\end{lstlisting}

String literals are concatenated (\texttt{"wor" "ld!"} becomes \texttt{"world!"}).
